\documentclass[a4paper,12pt]{article}
\usepackage[UTF8]{ctex}

\usepackage{enumerate}

\title{ML-Cheet-Sheet项目工作记录}

%%%%%%%%%%%%%%%%%%%%%%%%%%%%%%%%%%%%%%%%%%%%%%%%%%%%%%%%%%%%%%%%%%%%%%%%%%%%%%%
%%%%%%%%%%%%%%%%%%%%%%%%%%%%%%%%%%%%%%%%%%%%%%%%%%%%%%%%%%%%%%%%%%%%%%%%%%%%%%%
%%%%%%%%%%%%%%%%%%%%%%%%%%%%%%%%%%%%%%%%%%%%%%%%%%%%%%%%%%%%%%%%%%%%%%%%%%%%%%%
\begin{document}
\maketitle

\section{操作流程}
\subsection{项目文件结构}
\begin{itemize}
    \item 文件夹to-do：此文件夹中放置所有需要我们解释和整理的机器学习/深度学习概念图片。
    \item 文件夹done：此文件夹中放置所有已经完成解释和整理工作的概念图片。
    \item 文件夹illustration：此文件夹中放置在解释过程需要的额外图片。
    \item main.tex：项目的主文件，所有的概念解释均在这个latex文件中完成。
    \item work record.tex：项目进度记录文件。
\end{itemize}

\subsection{工作步骤}
\begin{enumerate}
    \item 进入项目后，请在to-do文件夹中寻找自己想要进行解释的概念，并将其拖拽到done文件夹。
    \item 进入main.tex文件，观察目录中是否有你想要解释的概念的正确分类，如果没有，请在文档的合适位置建立你需要的分类，分类的标题级别为section.
    \item 在正确的section中插入你想要解释的概念，这一步最好能注意一下同类概念有没有必要的先后顺序。
    \item 一个完整的概念解释包含一下几个部分：标题，即概念的名称，使用中文加英文的格式，使用subsection标题级别；概念图片；【概念解释】；【数学推导】；【概念对比】；【注意点】等。
    \item 结束的时候，将自己的工作记录在本文件的工作记录中，格式为姓名：完成解释的概念。
\end{enumerate}
%%%%%%%%%%%%%%%%%%%%%%%%%%%%%%%%%%%%%%%%%%%%%%%%%%%%%%%%%%%%%%%%%%%%%%%%%%%%%%%
%%%%%%%%%%%%%%%%%%%%%%%%%%%%%%%%%%%%%%%%%%%%%%%%%%%%%%%%%%%%%%%%%%%%%%%%%%%%%%%
%%%%%%%%%%%%%%%%%%%%%%%%%%%%%%%%%%%%%%%%%%%%%%%%%%%%%%%%%%%%%%%%%%%%%%%%%%%%%%%

\section{工作记录}
\subsection{2020/09/07 - 2020/09/13}
\begin{enumerate}
    \item 开始招募初始团队，通过郭金珊建立了四个人的小团队开始下一步工作的展开。
\end{enumerate}

\subsection{2020/09/14 - 2020/09/20}
\begin{enumerate}
    \item 第一次小组成员会议召开，尝试使用Github来管理项目版本和同步。后因为Github的Markdown版本全面支持数学公式的显示而放弃。
    \item 在overleaf平台建立项目，用于多人同步开发。
    \item 张海飞：误差Bias
   
\end{enumerate}

\subsection{2020/09/21 - 2020/09/27}
\begin{enumerate}
    \item 张海飞：AIC池化信息准则，混淆矩阵
    \item 郭金珊：方差，误差，完善偏差Bias，偏差-方差均衡
\end{enumerate}

\subsection{2020/09/28 - 2020/10/04}
\begin{enumerate}
    \item 郭金珊：损失函数VS风险函数VS目标函数，0-1损失，log对数损失
\end{enumerate}

\subsection{2020/10/05 - 2020/10/11}
\begin{enumerate}
    \item 郭金珊：学习，训练集VS验证集VS测试集，泛化能力，训练误差 VS 测试误差，假设空间，均方误差损失，绝对平均损失，Hinge损失，汉明损失，交叉熵损失，损失函数最小值，没有免费午餐定理，微调混淆矩阵以及引申的相关性能评估指标
    
    \item 张海飞：调整部分分类，
\end{enumerate}

\subsection{2020/10/12 - 2020/10/18}
\begin{enumerate}
    \item 张海飞： 样本矩阵，类别特征，One-hot编码，有序类别特征的编码，增加section特征工程 
\end{enumerate}

\subsection{2020/10/19 - 2020/10/25}
\begin{enumerate}
    \item 张海飞：贝叶斯错误率，欠采样，过采样 
\end{enumerate}

\subsection{2020/10/26 - 2020/11/01}
\begin{enumerate}
    \item 张海飞：监督学习VS非监督学习，分类，回归 
\end{enumerate}

\subsection{2020/11/02 - 2020/11/08}
\begin{enumerate}
    \item 张海飞： AUC，Bootstrapping
\end{enumerate}

\subsection{2020/11/16 - 2020/11/22}
\begin{enumerate}
    \item 郭金珊： 凹凸函数，R方决定系数
\end{enumerate}

\subsection{2020/11/23 - 2020/11/29}
\begin{enumerate}
    \item 张海飞： 维数灾难，特征选择
    \item 朱靖：完善1.3.1基本公式，1.3.2变量定义
\end{enumerate}

\subsection{2020/11/30 - 2020/12/06}
\begin{enumerate}
    \item 张海飞： 缺失值填充，使用k近邻进行缺失值填充
\end{enumerate}

\subsection{2020/12/07 - 2020/12/13}
\begin{enumerate}
    \item 张海飞： 基于树的算法，包含决策树，决策树回归
    \item 李慧：矩阵
    \item 完善点估计之矩估计法，以及大数定理
    \item 郭金珊：完善部分已写内容，欠拟合和过拟合，(非)参数模型，贝叶斯法则
\end{enumerate}

\subsection{2020/12/14 - 2020/12/20}
\begin{enumerate}
    \item 张海飞： Bagging，Random Forest，Boosting，AdaBoost
    \item 郭金珊： 决策规则，最大似然估计器，贝叶斯预测器
\end{enumerate}


\subsection{2020/12/14 - 2020/12/20}
\begin{enumerate}
    \item 张海飞： 特征重要性，热点对特征重要性的影响，决策边界
    \item 李慧：矩阵乘法
\end{enumerate}

\subsection{2021/01/11 - 2021/01/17}
\begin{enumerate}
    \item 朱靖 ：完善1.3.2随机变量及其分布，以及部分1.3.3多维随机变量
    \item 张海飞：逻辑回归，多元逻辑回归
\end{enumerate}


\subsection{2021/01/18 - 2021/01/24}
\begin{enumerate}
    \item 朱弘毅：导数
    \item 郭金珊：奥卡姆剃刀, 交叉验证法，完善补充部分已写概念
\end{enumerate}

\subsection{2021/01/25 - 2021/01/31}
\begin{enumerate}
    \item 朱弘毅：费马引理，三大微分中值定理，洛必达法则
\end{enumerate}

\subsection{2021/02/01 - 2021/02/07}
\begin{enumerate}
    \item
\end{enumerate}

\subsection{2021/02/08 - 2021/02/14}
\begin{enumerate}
    \item 朱弘毅：泰勒公式，函数单调性判断，渐近线求法，函数凹凸性判断，弧微分，曲率，曲率半径
    \item 朱靖 ：随机变量的数字特征,数理统计的基本概念
\end{enumerate}

\subsection{2021/02/15 - 2021/02/21}
\begin{enumerate}
    \item 朱弘毅：平面曲线的切线法线，反函数/隐函数运算，检查高等数学部分全部内容
    \item 张海飞：错误类别，集成学习方法，KNN
\end{enumerate}

\subsection{2021/03/08 - 2021/03/14}
\begin{enumerate}
    \item 张海飞：多项式回归，岭回归
\end{enumerate}

\subsection{2021/03/15 - 2021/03/21}
\begin{enumerate}
    \item 张海飞：Lasso回归
\end{enumerate}

\subsection{2021/03/22 - 2021/03/28}
\begin{enumerate}
    \item 张海飞：弹性网络
\end{enumerate}

\subsection{2021/06/24}
\begin{enumerate}
    \item 朱弘毅：删除高数部分中的所有图片
\end{enumerate}

\end{document}
